\section{Purpose, Scope, and Decision Rule}
\label{sec:plan-purpose}

This document records the execution decisions behind the research proposal
\emph{Five Business Days}. It separates project management from the academic
argument while preserving the feasibility constraints that determine what can
be claimed. The proposal evaluates the February 2024 acceleration of the
initial Schedule~13D filing deadline, with required trading days as the primary
exposure measure, derivative substitution as a second behavioural margin, and
the timing of takeover gains as a complementary measurement exercise.

The earlier 92-page liquidity-hump draft remains a separate working paper. It
is frozen after a bounded hygiene pass covering the enumerated cross-country
threshold claims, the $\kappa\to1$ narrative, welfare units, and two citation
attributions. A one-week Gaussian-noise falsification exercise closes the
remaining question about the old model's hump and belongs in that draft's
appendix, not in this project.

The new project is gate-first. Full model derivations, appendix machinery,
calibration, and price-dependent empirical work begin only after the access and
EDGAR gates below establish that the required inputs exist and that at least one
non-mechanical margin can be measured. A failed gate changes the scope; it does
not suspend the project indefinitely.

The March 2027 objective is a milestone presentation containing the exposure
measure, first-stage evidence, a model skeleton, and an initial analysis of
target gains. It is not a draft deadline. The full paper is scheduled for
summer 2027 and the autumn job market.

\section{Design Decisions Already Fixed}
\label{sec:plan-fixed}

\subsection{Target-level exposure}

The original design used a leave-one-out mean of an activist's past
accumulation. That measure is not viable as the main source of identification.
\textcite{polk2024} report 9{,}685 initial Schedule~13D filings by 5{,}863
entities from 2001--2022, only 1.65 filings per entity. The project's initial
300-filing sample contains 259 unique filer CIKs, 29 filers with at least two
filings, and three with at least three. Combining the historical measure with
activist fixed effects would leave only a few dozen repeat funds observed on
both sides of the reform. It would also select the prominent activists most
able to pre-position or use derivatives.

The primary exposure is therefore the target-level number of trading days
required to add five percentage points of ownership at a participation cap,
\[
\mathrm{RTD}_i
=\frac{0.05\times S_{\mathrm{out},i}}
       {\phi\times\mathrm{ADV}_i}.
\]
Activist history remains a secondary heterogeneity split for activists with at
least two pre-2024 non-corporate-action campaigns.

\subsection{Outcomes that cannot identify behaviour}

Filing-delay compression is required by the new rule and cannot by itself show
that activist behaviour changed. The two-business-day amendment deadline also
changes post-filing accumulation independently of the initial-filing deadline.
Neither measure can pass an empirical gate. The behavioural margins are stake
at filing, trigger-to-filing accumulation, campaign entry and target
composition, and Item~6 derivative use.

\subsection{Model and gain measurement}

The model is restricted to one private type, a quadratic post-trigger
accumulation cost, an exogenously garbled filing, a free-entry bidder cutoff,
and a bargaining share disciplined by one appropriability parameter. The model
does not add a second operational-versus-sale type, an endogenous signalling
game, the tender-game primitives, a within-window trading path, or priced
cash-settled derivatives. Ten to fourteen pages are feasible for the restricted
model; the discarded specification would require roughly eighteen to
twenty-two pages plus an appendix. Any new mechanism requires removing an
existing one.

The target-gain analysis uses a fixed event-time anchor. Calendar-window
decompositions are invalid here because the reform mechanically moves days from
the filing interval to the post-filing run-up. The main filing object is the
$[-1,+1]$ filing return as a share of total gain, holding the trigger-to-bid
horizon fixed. A mechanical placebo recuts pre-reform filings at five business
days, and a behavioural placebo uses campaigns that already filed within five
business days under the old rule.

\section{Data Inventory and Access Status}
\label{sec:plan-data}

The core estimation sample is 2022--2026. The SEC's published 2011--2021 tables
anchor the earlier period; a full 2010--2021 backfill begins only after
Item~5(c) extraction is validated on the core sample. Table~\ref{tab:plan-data}
records the build status as of 2 August 2026.

\begin{table}[htbp]
\centering
\scriptsize
\caption{Data and implementation inventory}
\label{tab:plan-data}
\begin{tabular}{@{}p{0.22\linewidth}p{0.10\linewidth}p{0.13\linewidth}p{0.17\linewidth}p{0.27\linewidth}@{}}
\toprule
Component & Status & Build estimate & Gate & Principal risk \\
\midrule
Schedule~13D and amendment enumeration & partial & 1--2 weeks & none &
Low to medium; in progress on 2 August 2026. \\
Item~5(c) transaction tables & none & 4--8+ weeks & pre-estimation &
High; heterogeneous tables require programmatic extraction and substantial
manual review. \\
Item~4 purpose classification & none & 3--5 weeks & validation &
Medium to high; require 95 per cent automated-versus-manual agreement on a
blinded 200-filing subset. \\
Trigger-date parser & partial & 2--3 weeks & pre-estimation &
Current parse rate is 0.64--0.68; hand-check against Item~5(c). \\
Post-December-2024 XML ingestion & partial & 1--2 weeks & none &
Mechanically easier but covers only about 18 months. \\
CRSP or WRDS access & unconfirmed & 1--2 weeks after access & week-1 decision &
Required for average daily volume and all return-based outcomes. \\
Amihud and average daily volume & none & 1--2 weeks after CRSP & access &
High; required for the primary exposure measure. \\
M\&A outcomes & none & 4--8 weeks by EDGAR route & week-2 source decision &
High; choose Bloomberg access or hand-collection from Forms~8-K and DEFM14A. \\
Merger-background chronologies & none & out of scope & none &
No tooling, corpus, or design note exists; exclude from the critical path. \\
\bottomrule
\end{tabular}
\end{table}

Two access decisions dominate the price-based work. CRSP or WRDS determines
whether average daily volume, Amihud illiquidity, and returns can be constructed.
The M\&A source determines whether the gain analysis can cover the full deal
universe or only a hand-collected acquired-target subsample. The EDGAR first
stage does not depend on either decision.

\section{Pre-Analysis Gates}
\label{sec:plan-gates}

\subsection*{Gate 0, week 1: advisor ratification}

Prepare a two-page decision memo and hold a supervision meeting. The memo states
what is being abandoned as the job-market-paper framing, what remains frozen as
a separate working paper, and what the new project will build. Every subsequent
gate is conditional on this decision.

\subsection*{Gate 1, weeks 1--2: access and source decisions}

Submit and calendar the CRSP or WRDS request. Complete the Bloomberg terminal
inventory for deal fields, premia, export limits, and academic terms of use.
Choose between a Bloomberg deal route and hand-collection from EDGAR. If price
data remain unavailable, activate the EDGAR-only configuration rather than
delaying the first stage.

In parallel, complete the bounded hygiene pass and Gaussian-noise check on the
frozen companion draft. Neither task may expand into a redesign of that paper.

\subsection*{Gate 2, weeks 2--4: pre-registered EDGAR analysis}

Pre-register minimum detectable effects using the SEC's Tables~3, 5, and 6,
before examining the realised estimates. The EDGAR analysis contains:
\begin{enumerate}
\item untrimmed filing-delay distributions split by corporate-action status;
\item stake-at-filing distributions;
\item entry and composition by a coarse, price-free proxy for required trading
days;
\item Item~6 cash-settled-derivative language; and
\item switching between Schedules~13G and 13D near five per cent ownership.
\end{enumerate}

Proceed to the full build if at least one non-mechanical margin moves by its
pre-registered detectable magnitude. If none does, reduce the model to the
accumulation corner and gain-accounting restriction, and reposition the paper
around a well-measured policy null and the timing of target gains. Mechanical
filing compliance cannot satisfy this rule.

\section{Measurement and Validation Requirements}
\label{sec:plan-validation}

\subsection{Filing delay and corporate-action status}

The current pipeline reports a post-reform median delay of 5.0 business days in
2024Q3--Q4, compared with 7.0 days in 2023Q2--Q3. It also reports 75.6 per cent
of post-reform filings within five business days, compared with 35.7 per cent
before the reform, but the post-reform calculation removes
delays above 60 business days without documentation. In the full set of 96
non-missing post-period observations, 68 filings, or 70.8 per cent, fall within
five business days. The pooled data also combine corporate-action and
non-corporate-action filings, although the SEC reports late-filing rates of
34 and 11 per cent for those groups \parencite{sec2023modernization}.

The production statistic must therefore be untrimmed, with the 60-day trim
reported only as a robustness check, and must separate filings by
corporate-action status. The split depends on the Item~5(c) history classifier.

\subsection{Exposure construction}

The six-month average-daily-volume window ends 60 trading days before the
trigger. Item~5(c) transaction tables must be used to diagnose residual
activist trading in that window. The exposure measure is then reconstructed
with a 120-trading-day lag. Covariate balance across exposure bins is reported
before any estimate. Specifications include activist-by-quarter,
liquidity-decile-by-quarter, and size-decile-by-quarter fixed effects.

\subsection{Inference battery}

The causal interpretation requires all of the following:
\begin{itemize}
\item at least 500 placebo reform dates drawn from 2015--2023, with the complete
specification re-estimated at each date;
\item at least 500 permutations of required trading days across activists while
holding the panel structure fixed;
\item quarterly exposure-bin event studies beginning in 2019 and a joint
pre-trend test; and
\item a non-US activist-block sample based on UK DTR5 and E.U.\ Transparency
Directive notifications, used only to absorb the global M\&A and interest-rate
cycle, not as a legal-threshold comparison.
\end{itemize}

The realised estimate's percentile is reported in both placebo distributions.
Failure of the pre-trend condition prevents a causal interpretation. The
October 2023 to February 2024 period is excluded from the main pre-period, with
additional exclusions around the March 2022 proposal and April 2023 reopening.

\subsection{Power limits}

The SEC reports that only 3 per cent of non-corporate-action campaigns, about
seven per year, had accumulated less than 90 per cent of their stake by the new
deadline; only 1 per cent, about one per year, had accumulated less than
75 per cent \parencite{sec2023modernization}. The expected 2024--2026 sample is
roughly 540--700 non-corporate-action campaigns. The sale-intent triple
interaction on bid hazard has a minimum detectable effect of 17--31 percentage
points on a base rate of about 17 percentage points, with roughly three bids in
the materially exposed post-reform cell. It remains descriptive and is reported
with confidence intervals only.

\section{Eight-Week Execution Schedule}
\label{sec:plan-schedule}

\begin{table}[htbp]
\centering
\small
\caption{Critical path}
\label{tab:critical-path}
\begin{tabular}{@{}p{0.14\linewidth}p{0.36\linewidth}p{0.40\linewidth}@{}}
\toprule
Timing & Deliverable & Exit condition \\
\midrule
Week 1 & Advisor memo; CRSP or WRDS request; Bloomberg inventory & Gate 0
ratified and access requests documented. \\
Weeks 1--2 & Source decisions; companion-draft hygiene; Gaussian-noise check &
Price-data and M\&A routes fixed; companion work frozen. \\
Weeks 2--4 & Pre-registered EDGAR analysis & Non-mechanical proceed-or-reduce
decision recorded against ex ante minimum detectable effects. \\
Weeks 5--8 & Trigger parser; average daily volume or EDGAR turnover proxy;
Gate-2 write-up & Inputs required for the primary exposure design validated. \\
Weeks 5--10 & Item~4 purpose codebook & Blinded 200-filing validation reaches
95 per cent agreement. \\
Weeks 9--12 & Acquired-target pilot, conditional on M\&A route & Initial
filing, run-up, and markup decomposition completed. \\
\bottomrule
\end{tabular}
\end{table}

The trigger-date parser and the exposure denominator are the critical path.
The Item~4 codebook may extend to weeks 8--10, and the acquired-target pilot may
extend to weeks 9--12. Weeks 5--8 commit to the critical inputs and Gate-2
write-up, not to completion of those overflow tasks.

\clearpage
\section{Fallbacks, Dissemination, and Publication Strategy}
\label{sec:plan-fallbacks}

If the schedule slips, reduce scope in this order:
\begin{enumerate}
\item keep merger-background chronologies out;
\item omit the signalling block;
\item replace the full deal-universe analysis with the hand-collected
acquired-target subsample;
\item drop the 2010--2021 backfill; and
\item use the EDGAR-only configuration covering stakes, timing, composition,
and derivatives.
\end{enumerate}

The first-stage facts and exposure design go to SSRN when the exposure measure
is fixed, with month 4--5 as the target. The novelty search is repeated at every
revision, as is the publication-status check for
\textcite{celentanolevine2025}, whose capitalisation interpretation must remain
properly attributed. The literature characterisations are tied to full-text
reading records and logged quotations; the one claim based only on an
author-posted summary remains explicitly qualified.

Publication planning is contingent on the evidence. A detectable response in
campaign activity, derivative use, or both defines the upside set; a joint
response is its strongest case. These outcomes support an RFS-or-above target.
A precisely estimated null, bounded by the pre-registered minimum
detectable effects, is planned for JCF or RF. These are planning targets, not
claims made in the research proposal.
